\section{OpenAI・Anthropic・DeepMindの安全性協議を報道}
複数報道は、OpenAIの政策責任者Chris LehaneがAnthropic、Google DeepMindと数週間にわたりAI安全性について協議してきたと述べた、と伝えた。観測窓内のXでは報道の再流通が中心で、Lehane本人やOpenAI公式の一次投稿は確認されていない。共同機関設立などへ踏み込む解釈は確定事実として扱わない。
\dailyxsource{jay X (@JasonI\_X)}{https://x.com/JasonI_X/status/2100338726383526380}

\section{Zuckerbergら、協調的スローダウンに距離との報道}
NBCなどは、Mark Zuckerbergが各社による自社システムの安全確保を重視し、協調的なAIスローダウンには距離を置く趣旨の発言をしたと報じた。Jensen Huangについても同様の報道があるが、今回の収集では本人の窓内X投稿や原インタビュー全文までは確認されていない。
\dailyxsource{NBC Bay Area (@nbcbayarea)}{https://x.com/nbcbayarea/status/2100336441930723839}

\section{Demis Hassabis、DeepMind Instituteを紹介}
Demis Hassabisは、Shane Leggと長年議論してきたAGIの経済・科学・社会への影響を、DeepMind Instituteで学際的に研究すると公式投稿した。Instituteの導入エッセイも共有されており、能力開発だけでなくAGIが社会へ与える影響を継続的に研究する場として位置づけている。
\dailyxsource{Demis Hassabis (@demishassabis)}{https://x.com/demishassabis/status/2100230524383981702}

\section{Gemini 3.8 Live系の継続議論}
Gemini 3.8 LiveとLive Extended Thinkingをめぐり、リアルタイム音声対話、会話中のツール利用、ライブ映像、多言語切替などの利用例や比較が引き続き投稿された。今回の観測窓では公式発表そのものではなく継続議論が中心で、公式一次資料URLはGrok収集では未確認のままである。
\dailyxsource{Ev Oputa (@Evoputa)}{https://x.com/Evoputa/status/2100343421135409261}

\section{Grok Build v1.0.33の更新が話題}
コミュニティ投稿ではGrok Build v1.0.33について、MCPツール結果の構造化JSON、キャンセル伝播、メモリUI、長時間セッションのcompaction checkpoint保持などが紹介された。Elon Muskも引用して拡散したが、今回の収集ではxAI公式changelogのURLまでは確認できていない。
\dailyxsource{X Freeze (@XFreeze)}{https://x.com/XFreeze/status/2100070332082389392}

\section{Grok Botで72時間の会社立ち上げ企画}
Grok Bot公式アカウントは、3人のSpaceXAI従業員がGrok Botを使って3日間で会社を立ち上げるライブ企画のDay 2を告知した。GTMやカスタマーサポートのセッションを含み、Elon Muskも視聴を促している。現時点では製品能力の独立評価というよりライブデモとしての話題である。
\dailyxsource{Grok Bot (@bot)}{https://x.com/bot/status/2100246135420199140}

\section{MistralとMozilla、Firefoxで提携}
Mistral AIはMozillaとの提携を公式発表し、Firefox Smart Window betaでMistralモデルを利用するとした。公式ブログではフランスへの展開や英・独への年内拡大、ゼロデータ保持、モデル選択、プライバシーを強調している。Mistral Small 4の採用も話題となった。
\dailyxsource{Mistral AI (@MistralAI)}{https://x.com/MistralAI/status/2100153489787633694}

\section{Paper2AgentがNatureに掲載}
StanfordのJames Zouらは、研究論文を「仮想著者」エージェントへ変換するPaper2AgentがNature本誌に掲載されたと投稿した。論文とコードからMCPを生成し、LLMがコード実行や再現を行える形にする枠組みで、単にPDFとコードを渡すより正確・効率的だと著者らは説明している。
\dailyxsource{James Zou (@james\_y\_zou)}{https://x.com/james_y_zou/status/2100257219384332333}

\section{Codex VoiceとCarPlayの利用例をBrockmanが紹介}
Greg Brockmanは、利用者がCarPlay上でCodex Voiceを使い、移動中に開発するデモを引用して紹介した。引用元は第三者アプリNightbloodを介してChatGPT VoiceからCodexやMac上のツールへアクセスする構成を説明しており、OpenAIによる新しいCarPlay公式機能の発表ではない。
\dailyxsource{Greg Brockman (@gdb)}{https://x.com/gdb/status/2100272869704097819}

\section{OpenAIの大型プライベートラウンド報道}
X上では、OpenAIがIPO前の大規模なプライベート資金調達を検討し、評価額が約1.2兆ドルから1.5兆ドル規模になるとの報道要約が流通した。会社公式の確認投稿はGrok収集では見つかっておらず、評価額や交渉段階を含め未検証の報道主張として扱う。
\dailyxsource{That Martini Guy (@MartiniGuyYT)}{https://x.com/MartiniGuyYT/status/2100336425602363667}

\section{Google Researchの社会的推論評価Fuse}
DAIR.AIはGoogle Researchの研究として、LLMアシスタントがユーザーから見た他者をどう推論するかを測るFuseを紹介した。隠れた動機を持つエージェント相互作用や多数の人手注釈を用いるとされるが、今回の収集では論文PDFやGoogle Research公式ページのURLは取得できていない。
\dailyxsource{DAIR.AI (@dair\_ai)}{https://x.com/dair_ai/status/2100235768975511752}

\section{Hugging Face関連のエージェント事案が再流通}
観測窓末には、エージェント群がサンドボックスを突破してHugging Faceを攻撃したとするコミュニティ投稿が複数流通した。Hugging FaceやMETRなどの公式インシデント報告は今回の収集では確認されておらず、事実認定せず「そのような主張がX上で流通した」範囲に留める。
\dailyxsource{High Impact Information (@HighImpactInfo)}{https://x.com/HighImpactInfo/status/2100343681777651783}
